\documentclass[11pt,twoside]{book}
\usepackage[paperwidth=6in,paperheight=9in,inner=0.8in,outer=0.7in,top=0.85in,bottom=0.8in,headsep=0.25in]{geometry}
\usepackage{fontspec}
\usepackage{polyglossia}
\setmainlanguage{english}
\setotherlanguage{arabic}
\setmainfont{Times New Roman}
\newfontfamily\arabicfont[Script=Arabic,Scale=1.25]{DecoType Naskh}
\usepackage{fancyhdr}
\pagestyle{fancy}
\fancyhf{}
\fancyhead[LE]{\small\thepage\hspace{1.2em}\textsc{Readings in Classical Arabic Prose}}
\fancyhead[RO]{\small\textsc{Passage 12\enspace The Grammarian and the Bedouin}\hspace{1.2em}\thepage}
\renewcommand{\headrulewidth}{0pt}
\setcounter{page}{129}
\linespread{1.15}
\begin{document}
\noindent{\small\itshape Level: Part II. New constructions: the ``five nouns'' (\textup{أبوك، أخوك}); broken plurals of the pattern \textup{فَعَالِين}; diptote proper names; the reproachful \textup{هَلَّا}.}

\bigskip
\begin{Arabic}
\noindent\textbf{١٢: النَّحْوِيُّ وَالأَعْرَابِيُّ}

\medskip
قِيلَ: خَرَجَ نَحْوِيٌّ مِنْ أَهْلِ الكُوفَةِ إِلَى البَادِيَةِ يَطْلُبُ الفَصَاحَةَ مِنْ أَفْوَاهِ الأَعْرَابِ، فَأَقَامَ فِيهِمْ أَيَّامًا يَسْمَعُ وَيَكْتُبُ. فَلَقِيَ ذَاتَ يَوْمٍ أَعْرَابِيًّا عَلَى مَاءٍ لِبَنِي تَمِيمٍ، فَسَلَّمَ عَلَيْهِ وَجَلَسَ إِلَيْهِ.

قَالَ النَّحْوِيُّ: «يَا أَخَا العَرَبِ، كَيْفَ تَقُولُ: جَاءَ أَخُوكَ؟» قَالَ: «جَاءَ أَخُوكَ.» قَالَ: «فَكَيْفَ تَقُولُ: رَأَيْتُ...؟» قَالَ: «رَأَيْتُ أَخَاكَ.» قَالَ: «فَمَا لَكَ لَا تَقُولُ: رَأَيْتُ أَخُوكَ، كَمَا قُلْتَ: جَاءَ أَخُوكَ؟» فَقَالَ الأَعْرَابِيُّ: «وَاللهِ مَا أَدْرِي مَا تَقُولُ، غَيْرَ أَنَّ لِسَانِي لَا يَجْرِي بِهِ.»

فَأَرَادَ النَّحْوِيُّ أَنْ يَمْتَحِنَهُ فِي الجَمْعِ، فَقَالَ: «فَكَيْفَ تَجْمَعُ دُكَّانًا؟» قَالَ: «دَكَاكِينُ.» قَالَ: «فَسِرْحَانًا؟» قَالَ: «سَرَاحِينُ.» قَالَ: «فَثُعْبَانًا؟» قَالَ: «ثَعَابِينُ.» فَفَرِحَ النَّحْوِيُّ وَظَنَّ أَنَّهُ قَدْ وَقَعَ عَلَى قَاعِدَةٍ، فَقَالَ: «فَعُثْمَانَ؟» قَالَ: «عُثْمَانُونَ.»

فَضَحِكَ النَّحْوِيُّ وَقَالَ: «هَلَّا قُلْتَ: عَثَامِينُ، كَمَا قُلْتَ: سَرَاحِينُ وَثَعَابِينُ؟» فَغَضِبَ الأَعْرَابِيُّ وَقَالَ: «أَتُرِيدُ أَنْ أَجْمَعَ رَجُلًا سَمَّاهُ أَبُوهُ كَمَا أَجْمَعُ ذِئْبًا أَوْ حَيَّةً؟ هَذَا كَلَامٌ لَا نَعْرِفُهُ فِي بِلَادِنَا، وَلَا يَقُولُهُ أَحَدٌ مِنْ قَوْمِي.»

فَانْصَرَفَ النَّحْوِيُّ إِلَى الكُوفَةِ، وَكَانَ يَقُولُ لِأَصْحَابِهِ بَعْدَ ذَلِكَ: «تَعَلَّمْتُ فِي البَادِيَةِ أَنَّ لِلْعَرَبِ أَذْوَاقًا لَا تَحْوِيهَا القَوَاعِدُ، وَأَنَّ الأَعْرَابِيَّ يَعْرِفُ مَا لَا يَعْرِفُ لِمَاذَا يَعْرِفُهُ.»
\end{Arabic}

\bigskip
\noindent{\small\textbf{Notes.}\enspace
\textup{البادية} the desert, home of the Bedouin, held to speak the purest Arabic.\enspace
\textup{دُكّان، دكاكين} shop.\enspace
\textup{سِرْحان، سراحين} wolf.\enspace
\textup{ثُعْبان، ثعابين} serpent.\enspace
\textup{عُثْمان} a man's name: diptote, and as a proper name it takes the sound plural.\enspace
\textup{هَلَّا} ``why didn't you\ldots?'', a reproach.}

\bigskip
\noindent{\small\textbf{Exercises.}\enspace (1) Find every one of the ``five nouns'' in the passage and give its case. (2) Why does the Bedouin refuse \textup{عثامين}? (3) Form the broken plural of \textup{بُسْتان} and \textup{سُلْطان}.}
\end{document}
